\chapter{线性结构专项：链表、栈、队列和单调结构}

线性结构专项题看上去都很短，但它们是面试中最容易暴露基本功的一类。数组和字符串通常考察下标、窗口和局部状态；链表考察节点关系和指针顺序；栈考察最近未解决状态；队列考察层次推进；单调栈和单调队列考察候选淘汰。很多候选人的代码错在一个小边界，但真正的原因不是粗心，而是没有说清楚“当前结构里保存的对象到底代表什么”。

链表题首先要把“值”和“节点”分开。数组里交换两个元素，很多时候只要交换值；链表里更重要的是节点之间的连接关系。删除节点、反转链表、合并链表、K 个一组翻转，都不是在移动一段连续内存，而是在改变若干个 \code{next} 指针。写链表题之前，建议先在纸上画三段：已经处理好的部分、当前正在处理的部分、尚未处理的部分。每次改指针前都问一句：原来的后继是否已经保存？新的前驱是谁？最终头节点从哪里返回？

删除倒数第 N 个节点的暴力想法是先遍历一遍得到长度 \code{len}，再删除第 \code{len - n + 1} 个节点。这个做法完全正确，复杂度也是 \complexity{O(n)}，但需要两次遍历。面试更常希望你展示快慢指针：让快指针先走 \code{n} 步，然后快慢一起走；当快指针到达尾部时，慢指针正好停在待删除节点的前驱。优化的本质不是少了一行代码，而是把“倒数位置”转换成“两个指针之间保持固定距离”的不变量。

\begin{lstlisting}[language=C++,caption={LeetCode 19 删除倒数第 N 个节点：虚拟头节点和快慢指针}]
struct ListNode {
    int value = 0;
    ListNode* next = nullptr;
};

ListNode* remove_nth_from_end(ListNode* head, int n)
{
    ListNode dummy;
    dummy.next = head;

    ListNode* fast = &dummy;
    ListNode* slow = &dummy;

    for (int step = 0; step < n; ++step) {
        fast = fast->next;
    }
    while (fast->next != nullptr) {
        fast = fast->next;
        slow = slow->next;
    }

    ListNode* removed = slow->next;
    slow->next = removed->next;
    return dummy.next;
}
\end{lstlisting}

虚拟头节点的价值在这里非常明显。如果删除的是原来的头节点，没有虚拟头，就要单独写一段分支；有了虚拟头，删除头节点和删除中间节点统一成“修改某个前驱节点的 \code{next}”。这类题的面试表达要明确：快慢指针之间相隔 \code{n} 个真实节点，循环结束时快指针停在最后一个真实节点，慢指针停在待删节点前驱。边界要测三个：只有一个节点，删除头节点，删除尾节点。

合并两个有序链表是链表和双指针结合的题。暴力方法可以把两个链表的值全部放进数组，排序后重新建链表；这样会丢掉链表题真正考察的节点重连能力，也多用了 \complexity{O(n)} 额外空间。更自然的做法是维护一个结果链表尾指针，每次从两个链表头部取较小节点接到尾部。由于两个输入链表本来有序，每一步的局部最小节点就是全局下一个节点。

\begin{lstlisting}[language=C++,caption={LeetCode 21 合并两个有序链表：尾指针接入节点}]
ListNode* merge_two_lists(ListNode* first, ListNode* second)
{
    ListNode dummy;
    ListNode* tail = &dummy;

    while (first != nullptr && second != nullptr) {
        if (first->value <= second->value) {
            tail->next = first;
            first = first->next;
        } else {
            tail->next = second;
            second = second->next;
        }
        tail = tail->next;
    }

    tail->next = first != nullptr ? first : second;
    return dummy.next;
}
\end{lstlisting}

这段代码没有创建新节点，只是复用原链表节点。因此如果题目要求不能修改输入，就要改成新建节点版本。稳定性也值得说明：当两个节点值相等时，代码优先取 \code{first}，这会保持第一条链表中相等元素的相对顺序。复杂度是 \complexity{O(m+n)}，空间是 \complexity{O(1)}。错误高发点是忘记移动 \code{tail}，或者循环结束后忘记把剩余链表整体接上。

K 个一组翻转链表比普通反转多了“分组边界”。暴力想法是把所有节点放进数组，按 K 个一组交换指针或值，再重新连接；这样虽然容易写，但牺牲了链表题的原地处理能力。原地做法要先检查当前组是否有 K 个节点，不足 K 个就保持原样；足够 K 个时，只翻转这一段，并把上一组的尾部接到新头，把新尾接到下一组开头。

\begin{lstlisting}[language=C++,caption={LeetCode 25 K 个一组翻转链表：先找边界再局部反转}]
ListNode* reverse_k_group(ListNode* head, int k)
{
    if (k <= 1 || head == nullptr) {
        return head;
    }

    ListNode dummy;
    dummy.next = head;
    ListNode* group_previous = &dummy;

    while (true) {
        ListNode* kth = group_previous;
        for (int step = 0; step < k && kth != nullptr; ++step) {
            kth = kth->next;
        }
        if (kth == nullptr) {
            break;
        }

        ListNode* group_next = kth->next;
        ListNode* previous = group_next;
        ListNode* current = group_previous->next;

        while (current != group_next) {
            ListNode* next = current->next;
            current->next = previous;
            previous = current;
            current = next;
        }

        ListNode* new_group_tail = group_previous->next;
        group_previous->next = kth;
        group_previous = new_group_tail;
    }

    return dummy.next;
}
\end{lstlisting}

这题最重要的是命名。没有 \code{group_previous}、\code{kth}、\code{group_next} 这三个边界名，代码会变成一堆难以验证的指针操作。反转循环里让 \code{previous} 初始等于 \code{group_next}，是为了反转完成后原组头节点自然指向下一组开头。这个技巧能减少额外连接分支。面试时不要急着写代码，先画出一组翻转前后：上一组尾、当前组头、当前组尾、下一组头，四个位置说清楚后，代码才有依据。

栈题的核心是“最近未解决”。有效括号中，最近出现的左括号必须最先被匹配；路径简化中，最近进入的目录会被 \code{..} 回退；表达式求值中，最近的操作数和运算符构成局部归约。栈的优化不是为了降低所有问题的复杂度，而是把嵌套结构或最近依赖关系变成线性扫描。

最小栈是一个很好的设计题。只用一个普通栈保存元素，\code{top} 和 \code{pop} 都是常数时间，但 \code{getMin} 需要扫描所有元素。优化方法是额外维护一个同步最小值栈。每压入一个值，同时压入“截至当前栈顶的最小值”；弹出时两个栈一起弹。这样任意时刻最小值栈栈顶就是当前全部元素的最小值。

\begin{lstlisting}[language=C++,caption={LeetCode 155 最小栈：数据栈和最小值栈同步}]
class MinStack {
public:
    void push(int value)
    {
        this->values_.push_back(value);
        if (this->minimums_.empty()) {
            this->minimums_.push_back(value);
        } else {
            this->minimums_.push_back(std::min(value, this->minimums_.back()));
        }
    }

    void pop()
    {
        this->values_.pop_back();
        this->minimums_.pop_back();
    }

    int top() const
    {
        return this->values_.back();
    }

    int get_min() const
    {
        return this->minimums_.back();
    }

private:
    std::vector<int> values_;
    std::vector<int> minimums_;
};
\end{lstlisting}

这段代码故意使用 \code{std::vector} 作为栈，因为 \code{vector} 支持连续存储、\code{back}、\code{push_back} 和 \code{pop_back}，足够表达栈语义。类成员访问使用 \code{this->}，避免成员和局部变量混淆。常见错误是只在新值更小时才压入最小值栈，但弹出重复最小值时就会错。例如依次压入 \code{2, 2, 3}，如果只记录一个 \code{2}，弹出一个 \code{2} 后最小值仍应是 \code{2}。同步栈写法用空间换清晰和可靠。

单调栈是“最近未解决”的强化版：栈里保存的不是所有历史元素，而是一组仍有可能成为答案的候选。每日温度中，栈保存还没找到更高温度的下标，且栈中温度从底到顶单调不增。新温度到来时，如果高于栈顶，就说明栈顶那一天的答案已经确定，因为当前天是它右侧第一个更高温度。

柱状图中最大的矩形更能体现单调栈的威力。暴力做法枚举每根柱子作为矩形高度，再向左右扩展找到第一个更矮柱子，最坏 \complexity{O(n^2)}。优化的关键是：当我们知道某根柱子左右两侧第一个更矮的位置时，以它为最低高度的最大矩形宽度就确定了。单调递增栈可以在一次扫描中确定“右侧第一个更矮”。

\begin{lstlisting}[language=C++,caption={LeetCode 84 柱状图最大矩形：单调递增栈}]
int largest_rectangle_area(const std::vector<int>& heights)
{
    std::vector<int> stack;
    stack.reserve(heights.size() + 1);
    int answer = 0;

    for (int i = 0; i <= static_cast<int>(heights.size()); ++i) {
        const int current_height =
            i == static_cast<int>(heights.size()) ? 0 : heights[i];

        while (!stack.empty() && heights[stack.back()] > current_height) {
            const int height = heights[stack.back()];
            stack.pop_back();
            const int left_less_index = stack.empty() ? -1 : stack.back();
            const int width = i - left_less_index - 1;
            answer = std::max(answer, height * width);
        }
        stack.push_back(i);
    }

    return answer;
}
\end{lstlisting}

这里末尾人为加入高度 0 的哨兵，目的是把栈里剩余柱子全部弹出结算。弹出某个下标时，当前下标 \code{i} 是它右侧第一个更矮的位置；弹出后新的栈顶是它左侧第一个更矮的位置。因此宽度是两者之间的柱子数量。这个解释比“套单调栈模板”重要得多。复杂度是 \complexity{O(n)}，因为每个下标最多入栈一次、出栈一次。易错点是比较符号：使用 \code{>} 或 \code{>=} 都能做，但对重复高度的归属不同，必须和宽度计算保持一致。

队列和双端队列负责维护“按时间进入、按窗口过期”的候选。滑动窗口最大值的暴力解对每个窗口扫描一遍最大值，复杂度 \complexity{O(nk)}。堆也能做，但需要处理过期下标。单调队列更贴合题目：队头永远是当前窗口最大值下标，队尾负责删除不可能再成为最大值的候选。

\begin{lstlisting}[language=C++,caption={LeetCode 239 滑动窗口最大值：队头保存当前最大值下标}]
std::vector<int> max_sliding_window(const std::vector<int>& nums, int k)
{
    std::vector<int> answer;
    if (nums.empty() || k <= 0) {
        return answer;
    }

    std::deque<int> deque;
    for (int right = 0; right < static_cast<int>(nums.size()); ++right) {
        while (!deque.empty() && nums[deque.back()] <= nums[right]) {
            deque.pop_back();
        }
        deque.push_back(right);

        const int left = right - k + 1;
        if (deque.front() < left) {
            deque.pop_front();
        }
        if (left >= 0) {
            answer.push_back(nums[deque.front()]);
        }
    }
    return answer;
}
\end{lstlisting}

为什么队尾较小元素可以删除？设队尾下标为 \code{j}，新下标为 \code{right}，如果 \code{nums[j] <= nums[right]}，那么在未来任何同时包含 \code{j} 和 \code{right} 的窗口里，\code{j} 都不可能比 \code{right} 更优；而且 \code{j} 更早过期。因此删除 \code{j} 不会影响任何未来答案。这就是候选淘汰的证明。单调队列的复盘要说清两个动作：入队前从队尾删掉被新元素支配的候选；出答案前从队头删掉已经过期的候选。

环形链表是快慢指针的另一条主线。暴力方法可以把走过的节点地址放进哈希集合，每到一个节点就检查是否见过；如果见过，说明有环；如果走到空指针，说明无环。这个方法直观，时间 \complexity{O(n)}，空间 \complexity{O(n)}。快慢指针把空间降到常数：慢指针每次走一步，快指针每次走两步。如果没有环，快指针会先到空；如果有环，快指针进入环后会不断追赶慢指针，最终相遇。

\begin{lstlisting}[language=C++,caption={LeetCode 141 环形链表：快慢指针检测相遇}]
bool has_cycle(ListNode* head)
{
    ListNode* slow = head;
    ListNode* fast = head;

    while (fast != nullptr && fast->next != nullptr) {
        slow = slow->next;
        fast = fast->next->next;
        if (slow == fast) {
            return true;
        }
    }

    return false;
}
\end{lstlisting}

这题的正确性可以用相对速度解释。进入环之后，快指针相对慢指针每轮多走一步；环长度有限，因此相对距离会按环长度取模不断变化，必然出现 0，也就是相遇。不要把这题说成“快指针总会追上慢指针”就结束，最好补一句“因为在环内相对距离每次加一，对环长取模”。如果题目进一步要求找到入环节点，做法是在相遇后让一个指针回到头节点，两个指针都每次走一步，再次相遇的位置就是入口。这来自公式：头到入口距离加上入口到相遇点距离，与快慢指针走过路程差之间存在整环倍数关系。面试里不一定要完整推公式，但要知道这不是经验结论。

LRU 缓存是线性结构和哈希表组合的设计题。题目要求 \code{get} 和 \code{put} 都是 \complexity{O(1)}，并且容量满时淘汰最近最少使用的键。单独用数组或链表，可以维护使用顺序，但查找 key 是 \complexity{O(n)}；单独用哈希表可以常数查找，但无法常数时间移动“最近使用”的位置。正确组合是：双向链表维护从最近到最久的顺序，哈希表从 key 映射到链表节点。

\begin{lstlisting}[language=C++,caption={LeetCode 146 LRU 缓存：哈希表定位节点，双向链表维护顺序}]
class LRUCache {
public:
    explicit LRUCache(int capacity)
        : capacity_(capacity)
    {
    }

    int get(int key)
    {
        const auto found = this->items_.find(key);
        if (found == this->items_.end()) {
            return -1;
        }
        this->touch(found->second);
        return found->second->second;
    }

    void put(int key, int value)
    {
        const auto found = this->items_.find(key);
        if (found != this->items_.end()) {
            found->second->second = value;
            this->touch(found->second);
            return;
        }

        this->entries_.emplace_front(key, value);
        this->items_[key] = this->entries_.begin();

        if (static_cast<int>(this->entries_.size()) > this->capacity_) {
            const int removed_key = this->entries_.back().first;
            this->items_.erase(removed_key);
            this->entries_.pop_back();
        }
    }

private:
    using Entry = std::pair<int, int>;
    using Iterator = std::list<Entry>::iterator;

    void touch(Iterator iterator)
    {
        this->entries_.splice(this->entries_.begin(), this->entries_, iterator);
    }

    int capacity_ = 0;
    std::list<Entry> entries_;
    std::unordered_map<int, Iterator> items_;
};
\end{lstlisting}

这段代码使用 \code{std::list}，不是因为链表在刷题中通常更快，而是因为 LRU 需要在已知节点迭代器的情况下常数时间移动节点到表头。 \code{splice} 不复制节点，只改变链表连接关系；哈希表保存的迭代器仍指向同一节点。这里 \code{std::list} 的节点稳定性正好匹配需求。反过来，如果用 \code{vector} 保存顺序，移动到表头会导致大量元素移动，还可能让保存的下标或迭代器失效。LRU 的复盘重点是“为什么必须两个结构一起用”：哈希表负责定位，双向链表负责顺序，任何一个单独使用都达不到题目复杂度。

接雨水可以用多个思路解决。暴力方法枚举每个位置，向左找最高柱子，向右找最高柱子，该位置能接的水是左右最高的较小值减当前高度；每个位置都向两边扫描，复杂度 \complexity{O(n^2)}。前后缀数组可以把左右最高值预处理出来，时间降到 \complexity{O(n)}，空间 \complexity{O(n)}。双指针进一步把空间降到 \complexity{O(1)}：左右两侧各维护当前最高柱子，哪边较低，就结算哪边。

\begin{lstlisting}[language=C++,caption={LeetCode 42 接雨水：双指针维护两侧最高柱}]
int trap(const std::vector<int>& height)
{
    int left = 0;
    int right = static_cast<int>(height.size()) - 1;
    int left_max = 0;
    int right_max = 0;
    int water = 0;

    while (left < right) {
        if (height[left] < height[right]) {
            left_max = std::max(left_max, height[left]);
            water += left_max - height[left];
            ++left;
        } else {
            right_max = std::max(right_max, height[right]);
            water += right_max - height[right];
            --right;
        }
    }

    return water;
}
\end{lstlisting}

为什么可以结算较低的一侧？假设 \code{height[left] < height[right]}，右侧至少存在一个高度大于左侧当前位置的柱子，因此当前位置水位的右边界不会低于 \code{height[left]}；它能接多少水只取决于左侧最高柱 \code{left_max}。如果左侧最高还不够高，就接不到水；如果够高，就能立刻计算。单调栈也能做接雨水，它在遇到右边界时弹出低洼底部并结算横向宽度。双指针更节省空间，单调栈更接近“找左右边界”的模型，面试中可以比较两者。

表达式求值考察栈对局部归约的建模。最简单的四则表达式如果没有括号，但有乘除优先级，可以扫描字符串，用栈保存已经归约到加减层面的项。遇到数字时，根据前一个运算符决定：加号就压入正数，减号压入负数，乘除就和栈顶立即合并。这样乘除不会等到最后才处理，加减则通过最后求和完成。

\begin{lstlisting}[language=C++,caption={LeetCode 227 基本计算器 II：栈处理乘除优先级}]
int calculate(const std::string& s)
{
    std::vector<int> stack;
    char operation = '+';
    int number = 0;

    for (int i = 0; i <= static_cast<int>(s.size()); ++i) {
        const char ch = i == static_cast<int>(s.size()) ? '+' : s[i];
        if (std::isdigit(static_cast<unsigned char>(ch))) {
            number = number * 10 + (ch - '0');
            continue;
        }
        if (ch == ' ') {
            continue;
        }

        if (operation == '+') {
            stack.push_back(number);
        } else if (operation == '-') {
            stack.push_back(-number);
        } else if (operation == '*') {
            stack.back() *= number;
        } else if (operation == '/') {
            stack.back() /= number;
        }

        operation = ch;
        number = 0;
    }

    return std::accumulate(stack.begin(), stack.end(), 0);
}
\end{lstlisting}

这里在循环末尾人为补一个加号，是为了触发最后一个数字的归约。表达式题的易错点非常多：多位数要持续累积；空格要跳过；除法在 C++ 中对整数向 0 截断，和题目语义是否一致要确认；如果有括号，就需要另一个层次的栈或显式状态机。不要一开始就写复杂解析器，先按题目支持的语法范围选择最小可行模型。面试表达时可以说：栈里保存的是已经处理完乘除优先级、等待最终加总的项。

线性结构题的实验训练不应只跑样例。链表题要构造空链表、单节点、删除头、删除尾、K 不整除长度；栈题要构造重复最小值、空弹出是否题目保证、嵌套和交叉括号；单调结构要构造严格递增、严格递减、全相等、窗口为 1、窗口等于数组长度。面试中如果能主动说出这些测试，说明你不仅知道算法，也知道边界来自哪里。

相交链表是链表“节点身份”意识的经典题。两个链表如果相交，意思不是某个节点的值相等，而是从某个节点开始共享同一批节点对象。暴力方法可以把第一条链表所有节点地址放入哈希集合，再扫描第二条链表，第一个出现在集合中的节点就是交点，复杂度 \complexity{O(m+n)}，空间 \complexity{O(m)}。双指针优化的思路更巧：指针 A 走完第一条链后转到第二条链，指针 B 走完第二条链后转到第一条链。若两链相交，它们会在交点相遇；若不相交，最终同时到达空指针。

\begin{lstlisting}[language=C++,caption={LeetCode 160 相交链表：双指针走完两条路径}]
ListNode* get_intersection_node(ListNode* first, ListNode* second)
{
    ListNode* a = first;
    ListNode* b = second;

    while (a != b) {
        a = a == nullptr ? second : a->next;
        b = b == nullptr ? first : b->next;
    }

    return a;
}
\end{lstlisting}

为什么这段代码成立？设第一条链表独有长度为 \code{x}，第二条独有长度为 \code{y}，共享尾部长度为 \code{z}。指针 A 走过的路径是 \code{x + z + y} 后到达交点，指针 B 走过的路径是 \code{y + z + x} 后到达交点，总长度相同。若没有共享尾部，两个指针都会在走完 \code{m+n} 后变成空指针。这里不能比较节点值，因为不同节点可以有相同值；也不能修改链表结构，除非题目允许。

回文链表可以用数组保存所有值后双指针判断，时间 \complexity{O(n)}、空间 \complexity{O(n)}。如果要求 \complexity{O(1)} 额外空间，就要找到中点、反转后半段、逐个比较，最后最好把链表恢复。这个题综合考察快慢指针和链表反转。快指针每次走两步，慢指针每次走一步；当快指针到达尾部时，慢指针在中点附近。奇数长度时，中间节点不影响回文判断，可以让后半段从慢指针之后开始。

\begin{lstlisting}[language=C++,caption={LeetCode 234 回文链表：反转后半段后比较}]
ListNode* reverse_list(ListNode* head)
{
    ListNode* previous = nullptr;
    ListNode* current = head;
    while (current != nullptr) {
        ListNode* next = current->next;
        current->next = previous;
        previous = current;
        current = next;
    }
    return previous;
}

bool is_palindrome_list(ListNode* head)
{
    if (head == nullptr || head->next == nullptr) {
        return true;
    }

    ListNode* slow = head;
    ListNode* fast = head;
    while (fast->next != nullptr && fast->next->next != nullptr) {
        slow = slow->next;
        fast = fast->next->next;
    }

    ListNode* second_half = reverse_list(slow->next);
    ListNode* restore_head = second_half;
    ListNode* first_half = head;
    bool ok = true;

    while (second_half != nullptr) {
        if (first_half->value != second_half->value) {
            ok = false;
            break;
        }
        first_half = first_half->next;
        second_half = second_half->next;
    }

    slow->next = reverse_list(restore_head);
    return ok;
}
\end{lstlisting}

这段代码刻意恢复链表，因为真实工程里“检查是否回文”不应该悄悄改变输入结构。力扣题可能不检查恢复，但教材要训练副作用意识。中点循环使用 \code{fast->next != nullptr && fast->next->next != nullptr}，能让 \code{slow} 停在前半段尾部。比较时只遍历后半段长度即可。易错点是反转后半段之前没有保存头节点，导致无法恢复。

链表排序可以用归并排序。数组排序可以随机访问中点，链表不行；链表找中点要用快慢指针，并从中点断开。暴力做法是把节点值放进数组排序再写回链表，这改变的是值，不是节点顺序；如果节点携带更多字段或外部引用，这种做法就不等价。链表归并排序的优势是稳定、时间 \complexity{O(n\log n)}，额外空间如果用递归会有栈成本；为了符合工程中避免深递归的习惯，可以使用自底向上的迭代归并。

\begin{lstlisting}[language=C++,caption={LeetCode 148 链表排序：自底向上迭代归并}]
ListNode* split_list(ListNode* head, int size)
{
    for (int i = 1; head != nullptr && i < size; ++i) {
        head = head->next;
    }
    if (head == nullptr) {
        return nullptr;
    }
    ListNode* second = head->next;
    head->next = nullptr;
    return second;
}

ListNode* merge_sorted_lists(ListNode* first, ListNode* second, ListNode** tail)
{
    ListNode dummy;
    ListNode* current = &dummy;
    while (first != nullptr && second != nullptr) {
        if (first->value <= second->value) {
            current->next = first;
            first = first->next;
        } else {
            current->next = second;
            second = second->next;
        }
        current = current->next;
    }
    current->next = first != nullptr ? first : second;
    while (current->next != nullptr) {
        current = current->next;
    }
    *tail = current;
    return dummy.next;
}

ListNode* sort_list(ListNode* head)
{
    int length = 0;
    for (ListNode* node = head; node != nullptr; node = node->next) {
        ++length;
    }

    ListNode dummy;
    dummy.next = head;
    for (int size = 1; size < length; size *= 2) {
        ListNode* previous = &dummy;
        ListNode* current = dummy.next;
        while (current != nullptr) {
            ListNode* first = current;
            ListNode* second = split_list(first, size);
            current = split_list(second, size);
            ListNode* merged_tail = nullptr;
            previous->next = merge_sorted_lists(first, second, &merged_tail);
            previous = merged_tail;
        }
    }
    return dummy.next;
}
\end{lstlisting}

自底向上归并的思想是先合并长度为 1 的有序段，再合并长度为 2、4、8 的段，直到段长覆盖整个链表。每次 \code{split_list} 都把一段切断，避免合并时越界串到后续段。代码较长，但它把递归栈换成了显式循环。面试中如果时间有限，可以先讲递归归并，再说明生产环境可改成迭代版本。这里的质量点不是炫技，而是知道链表排序的节点重连语义和递归深度风险。

删除排序链表中的重复元素也有两个版本。简单版保留一个重复值，只要当前节点和下一个节点值相等，就跳过下一个节点。进阶版删除所有重复值，只保留原链表中没有重复出现的值；这时必须使用虚拟头节点和前驱指针，因为重复段可能从头节点开始。先识别一整段相等值，再决定是保留还是跳过整段。

\begin{lstlisting}[language=C++,caption={LeetCode 82 删除排序链表中的重复元素 II：识别重复段后整体跳过}]
ListNode* delete_duplicates_all(ListNode* head)
{
    ListNode dummy;
    dummy.next = head;
    ListNode* previous = &dummy;
    ListNode* current = head;

    while (current != nullptr) {
        bool duplicated = false;
        while (current->next != nullptr &&
               current->value == current->next->value) {
            duplicated = true;
            current = current->next;
        }
        if (duplicated) {
            previous->next = current->next;
        } else {
            previous = previous->next;
        }
        current = current->next;
    }

    return dummy.next;
}
\end{lstlisting}

这题的重点是 \code{previous} 的含义：它始终指向结果链表的尾部，也就是最后一个确认保留的节点。遇到重复段时，\code{previous} 不动，只修改它的后继；遇到非重复节点时，\code{previous} 才前进。边界包括全是重复值、头部重复、尾部重复、没有重复。链表题里，前驱指针比当前指针更重要，因为删除动作发生在前驱的 \code{next} 上。

路径简化是栈在字符串解析中的应用。Unix 路径中，多个斜杠等价于一个斜杠，\code{.} 表示当前目录，\code{..} 表示回到上一级目录。暴力按字符处理容易被连续斜杠和边界搞乱。更清晰的做法是把路径按斜杠切分成段，栈保存有效目录名：普通目录入栈，\code{..} 弹出一个目录，空段和 \code{.} 忽略。最后用栈中目录重建规范路径。

\begin{lstlisting}[language=C++,caption={LeetCode 71 简化路径：栈保存有效目录段}]
std::string simplify_path(const std::string& path)
{
    std::vector<std::string> stack;
    int index = 0;
    while (index < static_cast<int>(path.size())) {
        while (index < static_cast<int>(path.size()) && path[index] == '/') {
            ++index;
        }
        int start = index;
        while (index < static_cast<int>(path.size()) && path[index] != '/') {
            ++index;
        }
        const std::string part = path.substr(start, index - start);
        if (part.empty() || part == ".") {
            continue;
        }
        if (part == "..") {
            if (!stack.empty()) {
                stack.pop_back();
            }
        } else {
            stack.push_back(part);
        }
    }

    if (stack.empty()) {
        return "/";
    }
    std::string result;
    for (const std::string& part : stack) {
        result += "/";
        result += part;
    }
    return result;
}
\end{lstlisting}

这里的“栈”不是必须用 \code{std::stack}，用 \code{vector} 更方便遍历重建路径。容器选择要服务于操作：我们需要尾部压入弹出，也需要最后从头到尾遍历，所以 \code{vector} 很合适。若路径很长，频繁字符串拼接可以用 \code{ostringstream} 或先估算长度，但刷题中清晰优先。易错点是根目录下遇到 \code{..} 不能继续向上，只能保持根。

括号相关题还会扩展为“最长有效括号”。暴力枚举所有子串再判断是否合法，复杂度高。栈做法保存下标，而不是保存字符。栈底保存最近一个无法匹配的位置，作为当前有效段的左边界。遇到左括号入栈；遇到右括号先弹出一个左括号，若弹出后栈空，说明当前右括号无法匹配，把它的下标作为新边界；否则当前有效长度是当前下标减栈顶下标。

\begin{lstlisting}[language=C++,caption={LeetCode 32 最长有效括号：栈保存边界下标}]
int longest_valid_parentheses(const std::string& s)
{
    std::vector<int> stack;
    stack.push_back(-1);
    int answer = 0;

    for (int i = 0; i < static_cast<int>(s.size()); ++i) {
        if (s[i] == '(') {
            stack.push_back(i);
            continue;
        }

        stack.pop_back();
        if (stack.empty()) {
            stack.push_back(i);
        } else {
            answer = std::max(answer, i - stack.back());
        }
    }

    return answer;
}
\end{lstlisting}

这个题中 \code{-1} 是哨兵边界，表示在字符串开始之前有一个不可跨越位置。栈顶在计算长度时不是当前有效段的第一个字符，而是当前有效段左侧的边界。这个语义非常重要。若只保存括号字符，就无法计算长度；若只保存左括号数量，就无法处理中间断裂。单调栈、括号栈、路径栈都在提醒同一件事：栈里保存的不是“看起来像栈的东西”，而是后续计算真正需要的历史状态。

队列题的进阶点是层次和边界。普通 BFS 中队列保存待扩展节点；若需要输出每层，就在每轮开始记录队列大小；若需要最短步数，就把层数和状态一起维护，或每处理一层步数加一。双端队列出现在两类场景：一类是单调队列维护窗口候选，另一类是 0-1 BFS 中根据边权把状态放队首或队尾。它们都不是“队列模板”，而是对候选优先级的精确表达。

\begin{labbox}
线性结构实验：第一，用同一组链表节点分别测试删除倒数节点、回文判断、排序，检查函数是否意外破坏输入。第二，把最长有效括号的栈内容逐步打印出来，观察栈顶为何是边界下标。第三，分别用数组扫描、单调队列解决滑动窗口最大值，统计窗口大小增大时的运行差异。第四，给路径简化加入连续斜杠、根目录回退、隐藏文件名等测试，确认规则没有被字符串细节干扰。
\end{labbox}

\begin{principlebox}{线性结构题的统一问题}
链表问“哪条边会断，断之前有没有保存”；栈问“哪个历史状态最近且必须先解决”；队列问“当前层或当前窗口的边界在哪里”；单调结构问“哪个候选被新元素支配，为什么永远不可能再成为答案”。只要这四个问题能回答，线性结构题就不再是零散模板。
\end{principlebox}
